\Article{8}
We now wish to illustrate our general solution with a few examples. These will best bring out both how it is applied and the nature of some further circumstances to be considered in applying it.

First let both surfaces be planes, in which case we can set
\[
\begin{aligned}
x&=t,&y&=u,&z&=0\\
X&=T,&Y&=U,&Z&=0
\end{aligned}
\]
\noindent The differential equation
\[
\omega=\mathrm{d}t^2+\mathrm{d}u^2=0
\]
\noindent here gives the two integrals
\[
t+iu=\text{Const.},\quad t-iu=\text{Const.}
\]
\noindent and similarly the two integrals of the equation \(\Omega=\mathrm{d}T^2+\mathrm{d}U^2=0\) are
\[
T+iU=\text{Const.},\quad T-iU=\text{Const.}
\]
\noindent The two general solutions of the problem are therefore
\[
\begin{aligned}
\mathrm{I.}\quad&T+iU=f(t+iu),\quad T-iU=f'(t-iu)\\
\mathrm{II.}\quad&T+iU=f(t-iu),\quad T-iU=f'(t+iu).
\end{aligned}
\]
\par\indent\SourceLeaf{205}%
We can also express this result as follows. With the symbol \(f\) denoting an arbitrary function, we take the real part of \(f(x+iy)\) for \(X\) and the imaginary part, with the factor \(i\) omitted, for either \(Y\) or \(-Y\).

If we use the symbols \(\varphi,\varphi'\) in the sense of Article 7 and put
\[
\varphi(x+iy)=\xi+i\eta,\quad\varphi'(x-iy)=\xi-i\eta
\]
\noindent where \(\xi\) and \(\eta\) will evidently be real functions of \(x\) and \(y\), the \textit{first} solution gives
\[
\begin{aligned}
\mathrm{d}X+i\,\mathrm{d}Y&=(\xi+i\eta)(\mathrm{d}x+i\,\mathrm{d}y)\\
\mathrm{d}X-i\,\mathrm{d}Y&=(\xi-i\eta)(\mathrm{d}x-i\,\mathrm{d}y)
\end{aligned}
\]
\noindent and hence
\[
\begin{aligned}
\mathrm{d}X&=\xi\,\mathrm{d}x-\eta\,\mathrm{d}y\\
\mathrm{d}Y&=\eta\,\mathrm{d}x+\xi\,\mathrm{d}y
\end{aligned}
\]
\noindent If we now set
\[
\begin{aligned}
\xi&=\sigma\mathord{.}\cos\gamma,&\eta&=\sigma\mathord{.}\sin\gamma\\
\mathrm{d}x&=\mathrm{d}s\mathord{.}\cos g,&\mathrm{d}y&=\mathrm{d}s\mathord{.}\sin g\\
\mathrm{d}X&=\mathrm{d}S\mathord{.}\cos G,&\mathrm{d}Y&=\mathrm{d}S\mathord{.}\sin G
\end{aligned}
\]
\noindent so that \(\mathrm{d}s\) denotes a linear element in the first plane, \(g\) its inclination to the abscissa axis, \(\mathrm{d}S\) the corresponding linear element in the second plane, and \(G\) its inclination to the abscissa axis, the preceding equations give
\[
\begin{aligned}
\mathrm{d}S\mathord{.}\cos G&=\sigma\mathord{.}\mathrm{d}s\mathord{.}\cos(g+\gamma)\\
\mathrm{d}S\mathord{.}\sin G&=\sigma\mathord{.}\mathrm{d}s\mathord{.}\sin(g+\gamma)
\end{aligned}
\]
\noindent and hence, if we. as we may, take \(\sigma\) to be positive,
\[
\mathrm{d}S=\sigma\mathord{.}\mathrm{d}s,\quad G=g+\gamma
\]
\noindent We thus see (in agreement with Article 7) that \(\sigma\) represents the ratio of magnification of the element \(\mathrm{d}s\) in its representation \(\mathrm{d}S\), and. as it should be, is independent of \(g\). Similarly, the independence of the angle \(\gamma\) from \(g\) shows that all linear elements issuing from a point in the first plane
\SourceLeaf{206}
are represented by elements in the second plane that form the same angles with one another as the original elements and, we may add, \textit{in the same sense}.

If we choose a linear function for \(f\), so that \(f\upsilon=A+B\upsilon\), with constant coefficients of the form
\[
A=a+bi,\quad B=c+ei
\]
\noindent then
\[
\varphi\upsilon=B=c+ei
\]
\noindent and therefore
\[
\sigma=\BareRoot{(cc+ee)},\quad\gamma=\mathrm{Arc.}\,\mathrm{tang}\frac{e}{c}
\]
\noindent The ratio of magnification is therefore constant at every point, and the representation is similar in every respect to what it represents.

For any other function \(f\) the ratio of magnification will vary, as is easy to prove, and similarity can therefore hold only in the smallest parts.

If the positions in the representation corresponding to a specified number of given points in the first plane are prescribed, we can readily use the usual method of interpolation to find the simplest algebraic function \(f\) satisfying this condition. Denote the values of \(x+iy\) at the given points by \(a,b,c\) and so on, and the corresponding values of \(X+iY\) by \(A,B,C\) and so on. We must then set
\[
\begin{aligned}
f\upsilon={}&\frac{(\upsilon-b)(\upsilon-c).\ .\ .}{(a-b)(a-c).\ .\ .}\mathord{.}A
+\frac{(\upsilon-a)(\upsilon-c).\ .\ .}{(b-a)(b-c).\ .\ .}\mathord{.}B\\
&+\frac{(\upsilon-a)(\upsilon-b).\ .\ .}{(c-a)(c-b).\ .\ .}\mathord{.}C+\text{etc.}
\end{aligned}
\]
\noindent which is an algebraic function of \(\upsilon\) whose degree is one less than the number of prescribed points. With two points the function becomes linear, so there is perfect similarity.

This method has a useful application in geodesy. If the correct positions of a number of points are known, the method can improve a map based on measurements of middling quality that is good in its finer details but somewhat distorted as a whole. It goes without saying, however, that such a transformation must not extend far beyond the region those points enclose.

If we carry out the \textit{second} solution in the same way, we find that the only difference is that the similarity is reversed. All elements in the representation do form angles with one another equal to those in what is represented, but in the opposite sense, so that
\SourceLeaf{207}
what lies on the right there is on the left here. This difference, however, is inessential and disappears if in one of the planes we take as the lower side the side we had previously regarded as upper. This last remark applies whenever one of the two surfaces is a plane, so in the following examples of this kind we can confine ourselves to the first solution.

